\centerline{\bf \Huge Математические пострелушки.}

\begin{flushright}
        \scriptsize{Нужно сперва ввязаться в бой, а там видно будет}
\end{flushright}


\section{Оружейная 1}

\problem{1;1} В одном тиктот ролике каждый участник является либо Шиииш, либо Шуш. После того, как один из Шииииш решил стать Шуш, тех и других в фирме стало поровну. Затем еще три Шиииш решили стать Шуш, и тогда Шуш стало вдвое больше, чем Шиииш. Сколько участников в этой фирме?

\problem{2;1} 
Планктон заменил в примере на умножение двузначных чисел цифры буквами, в надежде узнать формулу крабсбургера: одинаковые – одинаковыми, разные – разными. У неё получилось $GO * LD = OOO$ Какая была итоговая формула?

\problem{3;1} 
Найдите восемь последовательных целых чисел, сумма первых пяти из которых равна сумме остальных трёх.

\problem{ы} У овец и кур вместе 36 голов и 100 ног. Сколько овец? 

\problem{4;2} Найдите все трехзначные числа, которые в 12 раз больше суммы своих цифр.

\problem{5;2} Подберите числа x, y и z так, чтобы выполнялось равенство $28x + 30y + 31z = 365.$




\newpage

\section{Оружейная 2}

\problem{1;1} 
 Подойдя к незнакомому  кораблю с одним подъемом и думая, что на каждой палубе по шесть лежанок, Ророноа Зоро решил, что нужная ему лежанка находится на четвертой палубе. Поднявшись на четвертую палубу, Ророноа Зоро обнаружил, что нужная ему лежанка действительно находится там, несмотря на то, что на каждой палубе по семь лежанок. Каким мог быть номер лежанки, в которую шел Зоро?


\problem{2;1}
Расставьте числа 1, 2, 3, ..., 10 по кругу так, чтобы разность любых двух соседей была равна 2 или 3.



\problem{3;2}
У скольких трёхзначных чисел цифра сотен больше цифры десятков?

\problem{4;2}
 Проводится следствие по делу об украденном борциге. Подозреваемых трое: Айта, Бембя и Данзан. На суде Данзан заявил, что борциг украл Бембя. Айта и Бембя тоже дали показания, но что они сказали, никто не запомнил, а все записи пропали. В ходе судебного заседания выяснилось, что борциг украл лишь один из подсудимых, и что только он дал правдивые показания. Так кто украл борциге?
 
\problem{5;2} На острове живут рыцари, которые всегда говорят правду, и лжецы, которые всегда лгут. Путник встретил троих островитян и спросил каждого из них: "Сколько рыцарей среди твоих спутников?". Первый ответил: "Ни одного". Второй сказал: "Один". Что сказал третий?

\problem{ыыы} Сумма двух натуральных чисел равна 1244. Если в конце первого приписать 3, а в конце второго отбросить 2, то числа окажутся равными. Найти эти числа.  

\newpage
\section{Оружейная 3}


\problem{ыы} 
В фотоателье залетели 50 птиц — 18 скворцов, 17 трясогузок и 15 дятлов. Каждый раз, как только фотограф щёлкнет затвором фотоаппарата, какая-то одна из птичек улетит (насовсем). Какое наибольшее число кадров сможет сделать фотограф, чтобы быть уверенным: у него останется не меньше 10 дятлов?

\problem{1;2} Имеется 19 гирек весом 1, 2, 3, ..., 19 г, из которых девять железных, девять бронзовых и одна золотая. Известно, что общий вес всех железных гирек на 90 г больше, чем общий вес бронзовых. Найдите вес золотой гирьки.

\problem{2;1} Один из пяти братьев испек маме пирог. Андрей сказал: "Это Витя или Толя". Витя сказал: "Это сделал не я и не Юра". Толя сказал: "Вы оба шутите". Дима сказал: "Нет, один из них сказал правду, а другой — нет". Юра сказал: "Нет, Дима, ты не прав". Мама знает, что трое из ее сыновей всегда говорят правду. Кто испек пирог?

\problem{3;1}
Сколько на шахматной доске имеется всевозможных прямоугольников, состоящих из четырех клеток?

\problem{4;2}
Фермер Вупсень вырастил огромную пшеницу и написал на ней огромное число, подряд натуральные числа от 1 до 500: 123...10111213...499500. Вредный кузнечик Пупсень съел немного пшеницы, и  от этого числа первые 500 цифр. Как вы думаете, с какой цифры начинается оставшееся число?

\problem{5;3} Окрашенный кубик с ребром 6 см распилили на кубики с ребром 1 см. Сколько будет кубиков с двумя окрашенными гранями?
